\clearpage
\section{이차와 삼차방정식}
\label{sec:quadratic-cubic-galois-groups}

\begin{learninggoal}
이차다항식의 갈루아군을 완전히 분류한다. 기약 삼차다항식의 갈루아군이 $A_3$ 또는 $S_3$임을 보이고, 판별식의 제곱근으로 두 경우를 구별한다.
\end{learninggoal}

\subsection{이차방정식}

$f=X^2+aX+b\in K[X]$이고 $\charac K\ne2$라 하자. 판별량
\[
  D=a^2-4b
\]
를 두면 근은
\[
  \frac{-a\pm\sqrt D}{2}
\]
이고 분해체는 $K(\sqrt D)$이다.

\begin{proposition}[이차다항식의 갈루아군]
$f=X^2+aX+b\in K[X]$가 분리이고 $\charac K\ne2$라 하자.
\begin{enumerate}[label=(\roman*)]
  \item $D$가 $K$의 제곱이면 $f$는 $K$에서 완전히 분해되고 갈루아군은 자명하다.
  \item $D$가 $K$의 제곱이 아니면 $f$는 기약이고
  \[
    \Gal(f/K)\cong C_2.
  \]
\end{enumerate}
\end{proposition}

\begin{proof}
근의 공식과 $K(\sqrt D)/K$의 차수를 사용하면 된다. 비자명한 경우의 유일한 비항등 자동동형은 $\sqrt D\mapsto-\sqrt D$이다.
\end{proof}

\subsection{삼차다항식을 우울형으로 바꾸기}

$\charac K\ne3$이고
\[
  g(X)=X^3+c_1X^2+c_2X+c_3
\]
라 하자. $X=Y-c_1/3$을 대입하면 이차항이 사라져
\[
  Y^3+aY+b
\]
꼴이 된다. 이 평행이동은 근들을 모두 같은 양만큼 옮기므로 분해체와 갈루아군을 바꾸지 않는다. 따라서 삼차다항식의 계산에서는
\[
  f=X^3+aX+b
\]
를 기본형으로 사용할 수 있다.

이제 $\charac K\ne2,3$이고 $f=X^3+aX+b$가 기약이라고 가정하자. 분해체를 $L$, 세 근을 $\alpha_1,\alpha_2,\alpha_3$라 하자. 기약성 때문에 갈루아군 $G$는 세 근 위에서 추이적으로 작용한다. $S_3$의 추이적 부분군은
\[
  A_3\cong C_3
  \qquad\text{또는}\qquad
  S_3
\]
뿐이므로 $|G|$는 $3$ 또는 $6$이다.

\begin{definition}[삼차다항식의 판별식]
\[
  \delta=(\alpha_1-\alpha_2)(\alpha_1-\alpha_3)(\alpha_2-\alpha_3)
\]
라 두고
\[
  \Delta_f:=\delta^2
\]
를 $f$의 판별식이라 한다. $f=X^3+aX+b$일 때
\[
  \Delta_f=-4a^3-27b^2.
\]
\end{definition}

마지막 공식의 일반적인 유도는 다음 장에서 종결식과 함께 다룬다. 여기서는 근과 계수의 관계를 전개하여 직접 확인할 수도 있다.

\begin{lemma}[교대곱의 변환]
$\sigma\in G\le S_3$가 근들에 유도하는 순열을 같은 기호로 쓰면
\[
  \sigma(\delta)=\sgn(\sigma)\delta.
\]
따라서 $\Delta_f=\delta^2$는 모든 자동동형에 의해 고정되어 $K$에 속한다.
\end{lemma}

\begin{proof}
근들의 순서를 바꾸면 반더몬드 곱은 순열의 부호만큼 변한다. 제곱하면 부호가 사라지므로 $\Delta_f$는 $G$ 전체가 고정한다. 유한 갈루아 확장의 전체 고정체가 $K$이므로 $\Delta_f\in K$이다.
\end{proof}

\begin{theorem}[기약 삼차다항식의 갈루아군 판정]
$\charac K\ne2,3$이고 $f=X^3+aX+b\in K[X]$가 기약이라 하자. 그러면
\[
  \Gal(f/K)\cong
  \begin{cases}
    A_3\cong C_3,& \Delta_f\in K^{\times2},\\
    S_3,& \Delta_f\notin K^{\times2}.
  \end{cases}
\]
여기서 $K^{\times2}$는 $K$의 영이 아닌 제곱들의 집합이다.
\end{theorem}

\begin{proof}
$G$는 $A_3$ 또는 $S_3$이다. 위 보조정리에서
\[
  G\subseteq A_3
  \quad\Longleftrightarrow\quad
  \sigma(\delta)=\delta\text{ for all }\sigma\in G
  \quad\Longleftrightarrow\quad
  \delta\in K.
\]
또한 $\delta^2=\Delta_f$이므로 $\delta\in K$인 것은 $\Delta_f$가 $K$에서 제곱인 것과 동치이다. $G$가 추이적이므로 $G\subseteq A_3$이면 $G=A_3$이고, 그렇지 않으면 홀순열을 포함하여 $G=S_3$이다.
\end{proof}

\begin{corollary}[판별식이 만드는 이차 중간체]
위 정리에서 $G\cong S_3$이면
\[
  K(\sqrt{\Delta_f})
\]
는 분해체 $L/K$ 안의 유일한 이차 중간체이고, $A_3$의 고정체이다.
\end{corollary}

\begin{proof}
$\delta^2=\Delta_f$이고 짝순열은 $\delta$를 고정하며 홀순열은 부호를 바꾼다. 따라서 $L^{A_3}=K(\delta)=K(\sqrt{\Delta_f})$. $S_3$의 지수 $2$ 부분군은 $A_3$ 하나뿐이다.
\end{proof}

\begin{example}[갈루아군이 대칭군인 삼차식]
\[
  f=X^3-X+1\in\Q[X]
\]
은 유리근을 갖지 않으므로 기약이다. 여기서 $a=-1$, $b=1$이므로
\[
  \Delta_f=-4(-1)^3-27=4-27=-23.
\]
$-23$은 $\Q$의 제곱이 아니므로
\[
  \Gal(f/\Q)\cong S_3.
\]
유일한 이차 중간체는 $\Q(\sqrt{-23})$이다.
\end{example}

\begin{example}[갈루아군이 순환군인 삼차식]
\[
  f=X^3-3X+1\in\Q[X]
\]
은 유리근을 갖지 않아 기약이고
\[
  \Delta_f=-4(-3)^3-27=108-27=81=9^2.
\]
따라서
\[
  \Gal(f/\Q)\cong A_3\cong C_3.
\]
이 경우 한 근을 첨가한 삼차체가 이미 분해체이다.
\end{example}

\begin{pitfall}
판별식만 계산하고 기약성 검사를 생략하면 안 된다. reducible cubic에서는 갈루아군이 자명군이나 $C_2$가 될 수도 있다. $A_3$ 또는 $S_3$라는 이분법은 분리 기약 삼차다항식에 대한 결론이다.
\end{pitfall}

\begin{checkpoint}
\begin{enumerate}[label=(\alph*)]
  \item $X^3-2$의 판별식을 구하고 갈루아군을 판정하라.
  \item 판별식이 제곱인 기약 삼차식의 분해체 차수가 왜 $3$인지 설명하라.
  \item $S_3$ 경우에 $K(\sqrt{\Delta_f})$가 유일한 이차 중간체인 이유를 군론으로 설명하라.
\end{enumerate}
\end{checkpoint}
